\begin{abstract}
We prove that all finite-order Hecke $L$-functions over
$\mathbb Q(\sqrt{-3})$ and all Dirichlet $L$-functions are zero-free in
the half-plane $\Re s>7/8$, with the principal pole at $s=1$ allowed.
In particular, the Riemann zeta function is zero-free in this half-plane,
proving the quasi-Riemann hypothesis.
\end{abstract}

\tableofcontents
\clearpage

\section{Introduction}\label{sec:introduction}

Let $F=\mathbb Q(\sqrt{-3})$, let $\mathcal O$ be its ring of integers,
and write $N\mathfrak a=|\mathcal O/\mathfrak a|$ for the norm of a
nonzero integral ideal. A finite-order Hecke character $\eta$ modulo
a nonzero integral ideal $\mathfrak f$ is a character of the ray class group
modulo $\mathfrak f$, extended by zero to ideals not coprime to $\mathfrak f$.
This agrees with the usual finite-order idelic definition: the complex
place contributes no nontrivial finite-order continuous character
\cite[Chapter VI, Section 1]{MilneCFT}. For $\Re s>1$, put
\[
 L_F(s,\eta)=\sum_{\mathfrak a\ne0}
       \frac{\eta(\mathfrak a)}{(N\mathfrak a)^s}
 =\prod_{\mathfrak p}
       \left(1-\eta(\mathfrak p)(N\mathfrak p)^{-s}\right)^{-1}.
\]
The same notation denotes its meromorphic continuation. For a
Dirichlet character $\chi$ modulo $q$, extended by zero on nonunits, we
likewise write
\[
 L(s,\chi)=\sum_{n\ge1}\frac{\chi(n)}{n^s}
 \qquad(\Re s>1)
\]
and then continue meromorphically; the character modulo $1$ gives the
Riemann zeta function $\zeta(s)$. The problem here is uniform exclusion:
to find a constant $\sigma_0<1$, independent of the character, its
conductor, and the height, such that these functions have no zeros in
$\Re s>\sigma_0$.

For $\zeta$ alone, the existence of a fixed $\sigma_0<1$ for which
$\zeta(s)\ne0$ in $\Re s>\sigma_0$ is called the quasi-Riemann
hypothesis \cite[Section 1]{BCSS}. It asks for one gap valid at every
height, not merely nonvanishing on $\Re s=1$. The corresponding uniform
formulation over all Dirichlet characters appears in
\cite[p.~288, Equation (1.4) and the following paragraph]{FG}.

The connection between these functions and prime distribution has a
long history. Dirichlet's 1837 proof of infinitude of primes in every
reduced arithmetic progression introduced the character $L$-series
that separate residue classes \cite{Dirichlet}. Riemann's 1859 memoir
related the zeros of $\zeta$ to the distribution of primes and formulated
the critical-line conjecture \cite{Riemann}. Hadamard and de la
Vall\'ee Poussin independently proved in 1896 that $\zeta$ has no zero
on $\Re s=1$, obtaining the prime number theorem
\cite{Hadamard,ValleePoussin}. Hecke subsequently developed the
character zeta functions and their analytic theory over number fields
\cite{Hecke}.

Classical zero-free regions for Hecke $L$-functions approach the line
$\Re s=1$ as the conductor or height increases and can allow one
simple real exceptional zero; a precise form for abelian extensions is given
by \cite[Theorem 3.1]{TZ}. Such regions do not give uniform exclusion in
a fixed half-plane. Zero-density estimates address a different question:
they bound the number of zeros to the right of a given vertical line.
For example, Guth and Maynard's large-value estimates improve zero-density
bounds for $\zeta$ \cite[Theorem 1.2]{GM}, but these zero-density bounds
do not exclude every zero.

\begin{theorem}\label{thm:main}
Every finite-order Hecke $L$-function over
$F=\mathbb Q(\sqrt{-3})$ has no zero in $\Re s>7/8$.
The same holds for every Dirichlet $L$-function, including $\zeta(s)$.
A pole at $s=1$ for a principal character is allowed.
\end{theorem}

In particular, Theorem~\ref{thm:main} resolves the quasi-Riemann
hypothesis affirmatively: the supremum of the real parts of the
nontrivial zeros of $\zeta$, namely its zeros in $0<\Re s<1$, is at most
$7/8$. The boundary $\Re s=7/8$ is not included, while any exceptional
real zero in $(7/8,1)$ is excluded. The theorem does not establish the
Riemann hypothesis or its generalized versions, which place nontrivial
zeros on $\Re s=1/2$. The Riemann hypothesis remains open \cite{Clay}.
